% Laget automatisk av elevark/lag_elevark.py fra sim/fritt-fall.html. Endringer her blir overskrevet.
\documentclass[11pt]{article}
\usepackage{elevark}
\begin{document}
\elevarkhode{Fysikk 1 · Rettlinjet bevegelse}{Fritt fall}{https://fysikklab.no/sim/fritt-fall.html}{Simuleringen}

Galileo lurte på om farten til en fallende gjenstand øker med strekningen den har falt, eller med tiden den har falt. Slipp eller kast en ball, ta et stroboskopbilde, og finn ut hvem som har rett.

\section*{Oppgaver}

\begin{oppgaveliste}
\oppgave \fase{gjett}{Gjett}Galileo hadde to hypoteser: 1) farten er proporsjonal med strekningen, $v = ks$, og 2) farten er proporsjonal med tiden, $v = kt$. Hvilken tror du stemmer? Slipp ballen fra 20~m og se på de to grafene for hypotesene.
\svar
\oppgave \fase{test}{Mål}Slipp ballen fra 20~m med 0,10~s mellom bildene. Bruk tabellen til å regne ut gjennomsnittsfarten mellom bildene. Hvor mye øker farten per sekund? Hva er da $g$?
\svar
\oppgave \fase{test}{Regn}Hvor lang tid bruker ballen på å falle 20~m? Hvor stor fart har den når den treffer bakken? Sjekk med simuleringen.
\svar
\oppgave \fase{test}{Regn}Ballen kastes rett opp med 10~m/s fra 20~m høyde. Hvor høyt kommer den? Hvor lang tid tar det før den treffer bakken, og med hvilken fart?
\svar
\oppgave \fase{forklar}{Forklar}Hva er farten og akselerasjonen i det høyeste punktet? Hvorfor er fartsgrafen en rett linje også når ballen snur?
\svar
\oppgave \fase{test}{Undersøk}Slipp ballen fra 20~m på Månen. Hvor mye lengre tid tar fallet enn på Jorda? Regn først.
\svar
\oppgave \fase{utfordring}{Utfordring}I oppgave 4: Hvor stor fart har ballen når den passerer starthøyden på vei ned? Forklar hvorfor, og test.
\svar
\end{oppgaveliste}

\end{document}
